После первых двух частей платёжная отрасль выглядит обозримой.
Модель движения денег, сообщений и обязательств описана; механизмы
доверия, подлинности и защиты данных разобраны; остаётся подставить
названия сетей. Именно на этом шаге начинаются дорогостоящие ошибки. Две транзакции выглядят одинаково на экране кассы,
проходят через похожие стадии авторизации и используют один и~тот же
пластик в~кошельке. При этом они живут по разным правилам,
с~разным набором документов, разной финальностью расчёта,
разной архитектурой маршрутизации и разными последствиями для банка,
продавца и процессинга.

Третья часть представляет эту неоднородность как набор сравнимых моделей.
В~третьей части книга переходит от общего словаря
к~сравнению трёх семейств систем: глобальные карточные
схемы; национальные и региональные карточные системы; A2A-системы
(account-based, см.~главу~\ref{ch:open-banking}), в~которых платёж
опирается на доступ к~счёту и~согласие клиента вместо карточного
идентификатора. К~каждому семейству задаются одни
и~те же вопросы: кто владеет правилами системы, какие документы
действительно главные, как разделены публичный свод
правил и~закрытая техническая спецификация, как достигается
окончательность расчёта и~что меняется для банка, продавца
и процессинга при выборе конкретной системы.

Часть устроена как последовательное сравнение нескольких типов
платёжных реальностей. Сначала идут
глобальные карточные схемы -- Visa и Mastercard. На них хорошо
видна классическая логика свода правил, авторизационной и клиринговой
систем, сертификации и сетевых сервисов. Затем рассматриваются
Мир, UnionPay и региональные карточные схемы СНГ -- варианты той же
карточной модели, на которых хорошо видны роль национальной
инфраструктуры, локального регулирования и политической экономики
платёжной сети.

Во второй половине части масштаб меняется ещё раз. СБП и~открытый
банкинг открывают мир, где оплата опирается на банковский счёт,
согласие на доступ и~другую семантику окончательности расчёта~---
без карточного идентификатора и~классического карточного жизненного
цикла с~отдельной системой диспутов. Эта смежная область разобрана
рядом с~карточными схемами, чтобы различия были видны на одном
уровне сравнения.

С~инженерной точки зрения это часть с~наибольшей долей
сравнительного анализа. Нужна дисциплина сравнения: к~любой системе
подходим с~одним и~тем же набором вопросов -- кто
управляет системой, где искать источник истины, как участники
подключаются и сертифицируются, что считается финальным расчётом
и~как распределяются риск, ответственность и коммерческие стимулы.

Глубина глав в~этой части намеренно асимметрична. СБП, Мир,
UnionPay и~открытый банкинг в~России разобраны
подробнее, чем Visa и~Mastercard, потому что книга адресована
русскоязычному инженеру, и~глубина пропорциональна релевантности
системы для его повседневной работы. Главы про Visa и~Mastercard дают опорную модель,
достаточную для сравнения и~для понимания, что сделано иначе
в~национальных и~A2A-системах. За~исчерпывающим описанием
международных схем читатель идёт в~их собственные своды правил
и~бюллетени.

Эта часть формирует у~читателя архитектурное и операционное мышление. Как только
сравнение систем становится конкретным, видно, почему одни системы
по-разному обрабатывают диспуты, иначе распределяют риск фрода
и~требуют иной сертификации. Видно и~другое: системы по-разному
поддерживают платежи по подписке и предъявляют разные требования
к~внутреннему процессингу. То, что в~предыдущей части выглядело
общим механизмом, здесь получает институциональную форму.

Из этого перехода вырастает следующая часть. После сравнения сводов
правил, расчётных моделей и~систем смещается фокус: теперь важно не
\emph{какая} система стоит за платёжным потоком, а~\emph{что
происходит с~этим потоком в~проде}. Этой стороне работы посвящена
часть~\ref{part:operations}, где на первый план выйдут
фрод, комплаенс, подписки, сверка, диспуты и~управление отказами.
